\clearpage
\SourceLeaf{276}
\renewcommand{\BookRunningTitle}{Die Perioden der Klassen}
\pdfbookmark[0]{Beweis einiger Sätze über die Perioden der Klassen}{periodes-classes}
\begin{center}
\Large\bfseries BEWEIS EINIGER SÄTZE ÜBER\par\medskip
\large DIE PERIODEN DER KLASSEN BINÄRER QUADRATISCHER FORMEN.
\end{center}
\BookRule{.18\textwidth}

\pdfbookmark[1]{Satz I}{periodes-theoreme-I}
\textsc{Satz. I.}\quad\emph{Die Anzahl der (eigentlich primitiven) Klassen derselben Determinante, die, zur $P$-ten Potenz erhoben, wobei $P$ entweder eine Primzahl oder eine Primzahlpotenz $=p^\pi$ ist, die Hauptklasse $K$ ergeben, ist entweder gleich $1$ oder gleich einer Potenz derselben Primzahl $p$.}

\emph{Beweis.}\quad Es sei $(\mathfrak Q)$ die ganze Gruppe aller in Betracht kommenden Klassen und $n$ ihre Anzahl. Da die Hauptklasse $K$ notwendig in $(\mathfrak Q)$ enthalten ist, ist der Satz evident, wenn sie die einzige darin ist. Gibt es aber noch andere, so wird die Anzahl der in der Periode jeder von ihnen enthaltenen Klassen eine Potenz von $p$ sein; es sei eine von ihnen $A$, und wir nehmen an, ihre Periode $(\mathfrak A)$ enthalte $p^\alpha$ Klassen, die sämtlich in $(\mathfrak Q)$ enthalten sein werden. Erschöpfen nun die Klassen dieser Periode $(\mathfrak A)$ die Gruppe $(\mathfrak Q)$, so gilt $p^\alpha=n$, und der Satz ist bewiesen; andernfalls sei $B$ eine beliebige Klasse von $(\mathfrak Q)$, die nicht in $(\mathfrak A)$ enthalten ist, und wir nehmen an, ihre Periode werde so weit entwickelt, bis wir zu einer Klasse $bB$ gelangen, die zugleich unter den Klassen von $(\mathfrak A)$ vorkommt. Dies muss notwendig eintreten, weil zumindest die Hauptklasse dieser Periode und $(\mathfrak A)$ gemeinsam ist. Nehmen wir nun an, $bB$ sei die \emph{erste} Klasse in der Periode von $B$, die $(\mathfrak A)$ gemeinsam ist, oder $b$ sei so klein wie möglich, so behaupte ich

$1^0$.\quad Dass $b$ eine Potenz von $p$ sein wird. Denn offenbar erhalten wir, wenn wir $b=p^\beta h$, $bB=iA$ und $hk\equiv1\pmod{p^\pi}$ setzen (was möglich sein wird), $kbB=p^\beta hkB=p^\beta B=ikA$,
\SourceLeaf{277}
das heißt, auch $p^\beta B$ wird unter den Klassen von $(\mathfrak A)$ vorkommen, woraus folgt, dass $h=1$ und $b=p^\beta$ gilt.

$2^0$.\quad Dass, wenn wir die Klassen $K,B,2B\dotsfour(b-1)B$ mit $(\mathfrak B)$ bezeichnen, sämtliche Kompositionen einer Klasse von $(\mathfrak A)$ mit einer Klasse von $(\mathfrak B)$ zusammen $p^{\alpha+\beta}$ verschiedene Klassen ergeben. Denn nehmen wir $mA+nB=m'A+n'B$ und $n=n'$ an, so gilt notwendig $m=m'$; ist $n>n'$, so gilt $(n-n')B=(m'-m)A$, was unmöglich ist, wenn nicht $n=n'$ gilt.

$3^0$.\quad Dass diese $p^{\alpha+\beta}$ verschiedenen Klassen in $(\mathfrak Q)$ enthalten sein werden, was evident ist.

Erschöpfen nun diese $p^{\alpha+\beta}$ Klassen die Gruppe $(\mathfrak Q)$, so ist der Satz bewiesen; andernfalls wählen wir eine weitere Klasse von $(\mathfrak Q)$, die nicht unter ihnen enthalten ist, nämlich $C$; wir setzen ihre Periode so weit fort, bis wir zu einer Klasse gelangen, die bereits unter den aus $(\mathfrak A)$ und $(\mathfrak B)$ zusammengesetzten Klassen enthalten ist. Durch eine ähnliche Schlussweise wie zuvor\textsuperscript{'} werden wir beweisen, dass der Exponent dieser Klasse eine Potenz von $p$, $=p^\gamma$, sein muss und dass die Komposition der ersten $p^\gamma$ Klassen der Periode von $C$ mit den bereits gefundenen $p^{\alpha+\beta}$ Klassen $p^{\alpha+\beta+\gamma}$ verschiedene Klassen ergibt, die sämtlich in $(\mathfrak Q)$ enthalten sind. Erschöpfen diese Klassen die Gruppe $(\mathfrak Q)$ noch nicht, so behandeln wir in derselben Weise eine vierte Klasse $D$ usw.; und es ist evident, dass, da $(\mathfrak Q)$ aus endlich vielen Klassen besteht, auch diese Operationen enden werden und $n$ gleich einer Potenz von $p$ sein wird.\quad \mbox{Q. E. D.}

\par\bigskip
\pdfbookmark[1]{Satz II}{periodes-theoreme-II}
\textsc{Satz. II.}\quad\emph{Wird die Anzahl aller Klassen des Hauptgeschlechts durch $a^\alpha b^\beta c^\gamma$ usw. ausgedrückt, wobei $a,b,c$ verschiedene Primzahlen bezeichnen, so gibt es in diesem Geschlecht $a^\alpha,b^\beta,c^\gamma$ usw. Klassen, die, jeweils zur Potenz $a^\alpha,b^\beta,c^\gamma$ usw. erhoben, die Hauptklasse ergeben.}

\emph{Beweis}\quad Es seien $A,A',A''$ usw. sämtliche Klassen, die, zur Potenz $a^\alpha$ erhoben, $K$ ergeben, und $(\mathfrak A)$ ihre Gesamtheit; ebenso $B,B',B''$ usw. $(\mathfrak B)$, $C,C',C''$, $(\mathfrak C)$ usw. usw. Ich behaupte, dass aus der Komposition aller Klassen von $(\mathfrak A)$ mit allen Klassen von $(\mathfrak B)$ mit allen Klassen von $(\mathfrak C)$ usw. Klassen hervorgehen, die voneinander verschieden sind. Denn ist $A+B+C\dots=A'+B'+C'\dots$ usw., so erhalten wir, wenn wir $A-A'=A''$, $B-B'=B''$ usw. setzen,
\[
A''+B''+C''\text{ usw.}=K
\]
also, wenn wir zur Potenz $b^\beta c^\gamma$ usw. erheben, $(b^\beta c^\gamma\dots)A''=K$, woraus leicht
\SourceLeaf{278}
$A''=K$ und $A=A'$ folgt, und ebenso gilt $B=B'$, $C=C'$ usw. Es sei die Gesamtheit dieser Klassen $=(S)$. Ferner ist klar, dass alle diese Klassen zum Hauptgeschlecht gehören werden. Schließlich kann es im Hauptgeschlecht keine Klasse geben, die nicht in $(S)$ enthalten ist. Es sei\ .\ .\ .
\BookRule{.34\textwidth}

\par\bigskip
\pdfbookmark[1]{Bemerkung}{periodes-bemerkung}
\begin{center}\textsc{Bemerkung.}\end{center}
\begingroup\small
Dieses 1801 geschriebene Fragment bezieht sich auf Disq. Arithm., Artikel~306, \textsc{ix}. Das Wort \emph{dignité} wird hier in einem sonst nicht üblichen Sinn gebraucht.
\begin{flushright}\textsc{Stern}.\end{flushright}
\endgroup
\BookRule{.34\textwidth}
